% Generated draft: requires mathematical and source review.
\documentclass[11pt]{article}
\usepackage[margin=25mm]{geometry}
\usepackage{fontspec}
\setmainfont{Latin Modern Roman}
\usepackage{amsmath,amssymb,mathtools,unicode-math}
\setmathfont{Latin Modern Math}
\usepackage{enumitem,needspace,xurl,hyperref}
\setlength{\parindent}{0pt}
\setlength{\parskip}{5pt}
\setlength{\emergencystretch}{3em}
\newcommand{\N}{\mathbb{N}}
\newcommand{\Z}{\mathbb{Z}}
\newcommand{\Q}{\mathbb{Q}}
\newcommand{\R}{\mathbb{R}}
\newcommand{\C}{\mathbb{C}}
\newcommand{\F}{\mathbb{F}}
\newcommand{\A}{\mathbb{A}}
\newcommand{\PP}{\mathbb{P}}
\newcommand{\E}{\mathbb{E}}
\newcommand{\eps}{\varepsilon}
\newcommand{\dd}{\,\mathrm{d}}
\begin{document}
\section*{UnsolvedMath ID: 30005521}
\textbf{Source problem code:} catalogue entry
\subsection{Definitions and mathematical statement}
The supplied record asks about a three-stage Compton-camera model with relatively large, spatially unresolved detectors: the final detector measures photon energy but does not record the detected photon's spatial location. The unknown in the extracted forward expression is a function \(f\). Its domain, admissible function class, and support assumptions are not specified.

The damaged source expression has the schematic form
\[
P^{(2)}(x_2,y,E'')
=\int\!\int\!\int
C_3(x_2,y,E',E'')\,
\widetilde C_2(x_1,x_2,E')\,
f(s)C_1(s,x_1)\,
\dd S(x_1)\,\dd S(s)\,\dd E'.
\]
This display records the visible integrand and differential factors only; it does not specify a defined operator. The extraction contains apparent energy-bound tokens \(E''\) and \(E_0\), and apparent integration-set tokens, but their placement and interpretation cannot be established from the supplied text. The kernels, integration surfaces, surface measures, energy conventions, and meanings of \(x_2\) and \(y\) are not supplied. In particular, the displayed dependence on spatial variables does not identify which variables remain observable after spatial resolution is removed.

A conditional formulation of the intended inversion question is as follows. Once the physical geometry, kernels, integration regions, energy range, admissible class \(\mathcal F\) of unknown functions, and detector observation rule are specified, let \(\mathcal T_3\) denote the resulting three-stage forward operator. Let \(\mathcal R\) denote the specified observation map that retains the measured energy data and any other recorded quantities while aggregating over the unrecorded final detection position. Define
\[
\mathcal T_{\mathrm{unres}}=\mathcal R\circ\mathcal T_3.
\]
Determine whether there exists a reconstruction map
\[
\mathcal I:\mathcal T_{\mathrm{unres}}(\mathcal F)\longrightarrow\mathcal F
\quad\text{such that}\quad
\mathcal I\bigl(\mathcal T_{\mathrm{unres}}f\bigr)=f
\qquad\text{for every }f\in\mathcal F,
\]
and, if so, obtain an inversion procedure. Here \(\mathcal F\), \(\mathcal T_3\), and \(\mathcal R\) are placeholders for missing model specifications, not additional hypotheses asserted by the source. The supplied record therefore identifies an inversion problem but does not determine a unique mathematical operator or a fully specified theorem statement. It also does not state whether stability or a particular reconstruction topology is required.
\subsection{Short English statement}
Can an unknown source distribution be reconstructed from the three-stage Compton-camera measurements when the final detector records photon energy but not the detection position? The supplied record leaves the precise forward model, detector aggregation rule, and admissible source class unspecified, so the statement above is conditional on those missing definitions. The record labels the problem open; that label is supplied curation, not an independently verified conclusion.
\subsection{Sources}
\begin{itemize}
\item Local UnsolvedMath registry: supplied record with permanent ID \(30005521\), title ``Compton Transform Inversion with Spatially Unresolved Detectors,'' exact source problem code \texttt{catalogue entry}, in the Oberwolfach Reports Open Problems collection. The recorded proposed year is 2023; no individual proposer is identified.
\item Supplied report citation: \emph{Tomographic Inverse Problems: Mathematical Challenges and Novel Applications} (2023), Oberwolfach Reports, pp.~1105--1194, DOI 10.4171/owr/2023/21. \url{https://doi.org/10.4171/owr/2023/21}. The record describes Oberwolfach Reports as documenting mathematical workshops and their open-problem discussions; no specific contribution, author, or problem page is supplied.
\item Supplied literature source, labeled ``unresolved-detector study'': \url{https://doi.org/10.1002/pamm.202300281}.
\item Supplied literature source, labeled ``resolved Compton inversion'': \url{https://doi.org/10.1117/1.JMI.7.3.032504}.
\item Supplied literature source, labeled ``Oberwolfach problem report'': \url{https://doi.org/10.4171/owr/2023/21}.
\item Literature assessment dated exactly \texttt{2026-08-22}: the supplied curation labels the spatially unresolved-detector extension open, states that exact inversion is available for spatially resolved data, and describes unresolved variants as modeled rather than generally inverted. This is supplied dated curation derived from the OWR record, not independently verified here; its status and cited mathematical claims require independent checking.
\item The registry assigns the statement status \texttt{corrected\_verified} and reports that the official report was checked directly and the clean statement restored definitions while removing extraction artifacts. These are supplied provenance claims, not verification performed for this draft; the supplied text still lacks the definitions needed to specify the transform.
\end{itemize}
\end{document}
